% sections/experiments.tex -- experimental setup and results (E0--E8).
% Sources: SPEC.md sections 5-8, experiments/PLAN.md, experiments/ENVIRONMENT.md,
% workloads/README.md, engine/README.md. Measured numbers only via \Num macros
% (generated/numbers.tex), tables via generated/tab_*.tex, figures ../figures/F*.pdf.
\section{Experiments}\label{sec:experiments}

The experiments answer nine questions, E0 to E8 (\cref{tab:sweeps}): whether the
exact methods are exact and how often the rule of version~1 is wrong (E0); how
the cost of an update depends on the size of the complex and on the size of the
update (E1, E2); how the methods behave on genuine three-dimensional topology
changes, on non-manifold complexes and on surfaces (E3--E5); where locality fails
(E6); how much the rebuild policy matters (E7); and what the Morse order buys
(E8). \Cref{sec:exp-platform,sec:exp-methods,sec:exp-workloads,sec:exp-timing,sec:exp-correctness}
describe the setup and \crefrange{sec:e0}{sec:e8} the results. Every number,
table and figure in the results is produced by a script from the raw
per-update records of the runs; none is typed by hand.

\subsection{Implementation and platform}\label{sec:exp-platform}

\paragraph{Engine.}
All methods except \method{gudhi} run in one C++17 program and share its
complex store, which maps the sorted vertex tuple of a simplex to its facets and
cofaces and validates every atomic edit. Order keys are dense 32-bit integers; a
(re)build assigns $0,\dots,N-1$ in the chosen order, and every later insertion
takes the next unused key, as \cref{sec:drb} requires. A chain is a strictly
increasing array of keys, so its low is the last entry and a $\Z$ addition is a
linear merge. When short chains are added to a long working column, the column
switches to a max-heap with lazy cancellation, so a reduction step costs
$O(|R_p|+|V_p|)$ up to a logarithmic factor, as in \cref{prop:cost}. The
footprint index keeps the exact count $|\occ(\sigma)|$ and a lazily filtered list
of slots, compacted when it exceeds twice the count plus a constant. In Case~A
the engine chooses the cycle $V_z$ with the fewest cells. The collapse of
\cref{def:collapse} uses a first-in-first-out queue initialised in (dimension,
vertex tuple) order and visits facets in a fixed order, so the matching depends
only on the complex, not on the history of edits. Debug builds check
(I1)--(I3), the pivot table and the footprint index after every operation.

\paragraph{Platform.}
All runs use one cloud virtual machine with an Intel Xeon processor at
2.10\,GHz (4~cores, one thread per core; 48\,KiB L1 data and 2\,MiB L2 cache per
core, 260\,MiB shared L3 cache) and 15\,GiB of memory, running Ubuntu 24.04.5
LTS with Linux~6.18. The engine is compiled by g++~13.3.0 (CMake 3.28.3) with
\texttt{-O2 -DNDEBUG} and no machine-specific tuning. The GUDHI baseline and
oracle run under Python~3.13.16 with GUDHI~3.13.0 and numpy~2.5.3. Every method
runs single-threaded.

\subsection{Methods}\label{sec:exp-methods}

\Cref{tab:methods} lists the methods; the names are those used in the code, the
raw data and the figures.

\begin{table}[t]
  \centering
  \caption{Methods. \emph{Order}: the key order of the cells. \emph{$V$,
    $\occ$}: whether the change of basis and the footprint index are stored.
    \emph{Policy}: the rebuild policy of \cref{sec:drb}. All methods except
    \method{gudhi} run in the same engine and share its complex store.}
  \label{tab:methods}
  \small
  \begin{tabularx}{\linewidth}{@{}l l l l l >{\raggedright\arraybackslash}X@{}}
    \toprule
    Method & Order & $V$, $\occ$ & Rebuilds & Exact & Role \\
    \midrule
    \method{full-id}    & insertion         & no       & every update & yes & static recomputation with clearing, ranks only \\
    \method{full-morse} & Morse, recomputed & no       & every update & yes & static recomputation in the Morse order of a new collapse matching \\
    \method{gudhi}      & GUDHI's own       & ---      & every update & yes & external static baseline: GUDHI simplex tree, persistence over $\Z$ \\
    \method{drb-id}     & insertion         & yes      & never        & yes & ablation: neither Morse order nor rebuilds \\
    \method{drb-id-rb}  & insertion         & yes      & policy       & yes & ablation: rebuilds without the Morse order \\
    \method{mmhm}       & Morse             & yes      & policy       & yes & the method of \cref{sec:theory} \\
    \method{mmhm-norb}  & Morse of $K_0$    & yes      & never        & yes & ablation: no rebuilds \\
    \midrule
    \method{lcr}        & Morse of $K_0$    & $R$ only & never        & no  & update rule of version~1 (defect E1); error rate only \\
    \method{euler}      & ---               & ---      & ---          & on surfaces & counters of \cref{rem:surfaces}; surfaces only \\
    \bottomrule
  \end{tabularx}
\end{table}

\paragraph{Static baselines.}
\method{full-id} and \method{full-morse} recompute all ranks after every update
by the static reduction with clearing (\cref{sec:drb}), without storing~$V$:
\method{full-id} in insertion order, \method{full-morse} in the Morse order of a
collapse matching computed anew after every update. \method{full-morse} is the
``coreduction-only'' baseline of version~1, done exactly: by
\cref{cor:critical-work,cor:morse-reduction} its column additions all occur in
critical columns and, after the gradient-flow steps, are those of the reduction
of the Morse complex, but it pays for a collapse after every update. Both share
the store and the reduction code with the dynamic methods, so differences
reflect the algorithms rather than the implementations.
\method{gudhi} replays the same edits on a GUDHI simplex
tree~\cite{maria2014gudhi} through its Python interface and, after every update,
calls \texttt{compute\_persistence} with the arguments
\texttt{homology\_coeff\_field=2} and \texttt{persistence\_dim\_max=True}, and then
\texttt{betti\_numbers}; without the second argument GUDHI omits the top
dimension. GUDHI computes persistence pairs, a richer output, and is not designed
for this task; it serves as a reference point for a widely used library and as
an oracle.

\paragraph{Dynamic methods.}
\method{mmhm} is the method of \cref{sec:theory}: a dynamic reduced basis built
in the Morse order and rebuilt in a fresh Morse order under the rebuild policy of
\cref{sec:drb}, with $m=\infty$ and $\alpha=1$ unless stated otherwise (E7
varies both). Three ablations separate its ingredients. \method{drb-id} uses
insertion keys and never rebuilds. \method{mmhm-norb} uses the Morse order of
$K_0$ and never rebuilds. \method{drb-id-rb} rebuilds under the same policy and
parameters as \method{mmhm}, but its rebuild reruns the static reduction with
clearing in the current insertion order (the keys in use, renumbered) instead
of computing a new matching. So \method{mmhm} against \method{drb-id-rb}
isolates the order used at a rebuild, and \method{mmhm} against
\method{mmhm-norb} isolates rebuilding.

\paragraph{Version 1's rule and the surface counters.}
\method{lcr} implements the update rule of version~1 without basis tracking
(defect E1 of \cref{sec:erratum}): an inserted column is reduced against the
existing pivots, and a deleted column releases its pivot and is dropped. It
uses exactly the order of \method{mmhm-norb} and differs from it only in not
storing $V$ and $\occ$, so its errors are caused by E1 alone (E2 concerns the
Morse complex, which \method{lcr} does not form). It is run only to measure how
often it is wrong. \method{euler} maintains $\chi$ and the numbers of edges with one, two or
more cofaces in $O(1)$ per atomic edit and returns the Betti numbers of
\cref{rem:surfaces}. At the end of every update it checks that the complex has a
triangle, no tetrahedron, and no edge outside one or two triangles, and aborts
otherwise; connectivity and the absence of pinched vertices are assumed, not
checked.

\subsection{Workloads}\label{sec:exp-workloads}

Five generators produce the update sequences. Each is
deterministic given its parameters and seed, emits only valid atomic edits
(checked again by an independent validator), lists deletions top-down and
insertions bottom-up, and records its parameters, the sizes of all states and
the number of edits per update. For W1, W2, W3 and W5 the Betti numbers of every
state are known in closed form and are written next to the sequence. Distances
are graph distances in the base mesh.

\paragraph{W1: punctured surfaces.}
The base surface is the geodesic icosphere of frequency $f=2^L$, with $10f^2+2$
vertices, $30f^2$ edges and $20f^2$ triangles, or the $n\times m$ periodic grid
torus with every square split along one diagonal, with $nm$ vertices, $3nm$
edges and $2nm$ triangles. For a vertex $v$, the patch $P(v,k)$ is the set of
triangles incident to a vertex at distance at most $k-1$ from~$v$; we use $k=2$.
When $v$ and its neighbours have degree six, $P(v,2)$ has $24$ triangles and
opening it takes $61$ atomic edits; no vertex of our surfaces has a larger
degree, so no update has more edits. A patch around a
random vertex is used only if its closure is a disk and its vertices are at
distance more than one from the vertices of every existing hole. Opening a hole
deletes the triangles of the patch, then its edges, then the vertices left
without cofaces; closing a hole is the exact reverse of its opening. With $h$
holes, the generator opens a hole if $h=0$; otherwise it opens one with
probability $0.75$, $0.5$ or $0.25$ when $h$ is below, at or above the target
$3$, and closes one otherwise or when no valid patch is found. The Betti numbers
are $(1,\max(h-1,0),[h=0],0)$ for the sphere and $(1,2+\max(h-1,0),[h=0],0)$ for
the torus, where $[h=0]$ is $1$ if $h=0$ and $0$ otherwise.
W1 is the setting of \cref{rem:surfaces}, which E5 examines. It also contains
the phenomenon of \cref{prop:nonlocal}: when no hole is open, the surface is a
connected closed surface, so the deletion that opens the first hole costs every
DRB method $\Omega(n_2)$ operations by \cref{prop:nonlocal}(a), and the last
triangle inserted when the last hole closes costs $\Omega(n_2)$ by
\cref{prop:nonlocal}(b).

\paragraph{W2: refinement and coarsening in a 3-ball.}
The base complex is the Freudenthal (Kuhn) triangulation of a grid of
$N\times N\times N$ cubes, each split into six tetrahedra; it has $(N+1)^3$
vertices and about $26N^3$ simplices. A refinement applies $s$ stellar $1\to4$
subdivisions to $s$ tetrahedra collected breadth-first over face adjacency from
a uniformly random tetrahedron; tetrahedra created by earlier subdivisions may
be subdivided again. A coarsening undoes $s$ leaf subdivisions, those whose four
tetrahedra are all still present: a uniformly random one and the $s-1$ nearest
to it, measured by the Euclidean distance between the inserted vertices. A
subdivision and its inverse each consist of $16$ atomic edits, so $s_t=16s$.
Because coarsening can only undo subdivisions, $K_0$ already contains $8s$ live
subdivisions, and with $x$ live subdivisions the generator refines with
probability $\min\bigl(0.95,\max(0.05,\,0.5-0.5(x-8s)/(8s))\bigr)$, and always
when fewer than $s$ leaf subdivisions exist. Since $|K|=|K_{\mathrm{base}}|+14x$,
the size performs a random walk around its initial value instead of growing
linearly (compare experiment P2 of version~1, \cref{sec:erratum}). The Betti
numbers are $(1,0,0,0)$ after every update, but not within an update: by
\cref{thm:exact}(d) every atomic edit changes one Betti number by one.

\paragraph{W3: cavities in a 3-ball.}
The base complex is that of W2. A centre $v$ is drawn uniformly among the
interior vertices, and the ball $B(v,k)$ is the set of tetrahedra incident to a
vertex at distance at most $k-1$ from~$v$; we use $k=2$. Carving deletes the
tetrahedra of $B$ and then, top-down, every face left without cofaces. A ball is
accepted only if it is more than one ring away from every carved ball, its
boundary is a $2$-sphere, and it meets the boundary of the cube in no triangle (a
cavity: $\beta_2$ grows by one) or in one disk (a dent: no change); other
contacts, which could bore a tunnel, are rejected. Refilling is the exact
reverse of a carving. The number of carved balls follows the same walk as the
number of holes in W1, with target $4$. The Betti numbers are $(1,0,c,0)$, where
$c$ is the number of cavities.

\paragraph{W4: flag complexes of moving points.}
$n$ points are drawn uniformly in the unit square or the unit cube (``2D'' and
``3D'' in the tables and figures). A radius
$\varepsilon$ is calibrated on the initial points so that the
$\varepsilon$-neighbourhood graph has exactly $5n$ edges (mean degree $10$), and
$K$ is the flag complex of this graph truncated at dimension $D\in\{2,3\}$. An
update moves one uniformly random point by adding to each coordinate an
independent offset drawn uniformly from $[-\varepsilon/4,\varepsilon/4]$,
reflected at the border of the domain, and keeps its vertex identifier: it
deletes the closed star of the vertex top-down and inserts the vertex and its new
star bottom-up. The complexes are not manifolds and all Betti numbers vary; they
are not known in closed form and are checked against the oracles only.

\emph{Caveat: truncation.} In a flag complex truncated at dimension $D$, the
vertices of a clique of size $D+2$ span a $(D+1)$-simplex that is not in~$K$, so
its boundary is a $D$-cycle that is never filled. Hence $\beta_D$ is large: at
INIT, $\beta_D/n$ lies between \NumWFourBetaDRatioMin{} and
\NumWFourBetaDRatioMax{} in our W4 instances (\cref{tab:w4-betti}). This is genuine linear algebra on
a non-manifold complex, but $\beta_D$ is dominated by the truncation, not by the
shape of the point cloud. We use W4 as an algebraic stress test and draw no
conclusion about the topology of point clouds from it.

\paragraph{W5: non-locality witnesses.}
(a) The cycle graph $C_L$, whose updates toggle the edge $\{0,1\}$: odd-numbered
updates delete it and even-numbered ones insert it again. (b) The icosphere of
frequency $2^L$, whose updates toggle one triangle chosen from the seed in the
same way. By \cref{prop:nonlocal}, every toggle costs every DRB method
$\Omega(L)$, respectively $\Omega(n_2)$, operations; once the states alternate,
a toggle costs $O(L+\nnz(R)+\nnz(V))$, respectively $O(n_2+\nnz(R)+\nnz(V))$.
\method{euler} needs $O(1)$ per toggle on (b) and does not apply to (a), which
is a graph and not a surface.

\paragraph{Sweeps.}
\Cref{tab:sweeps} lists the workloads, parameters and methods of each
experiment. Every timed configuration ran with \NumSeeds{} seeds.

\begin{table}[t]
  \centering
  \caption{Experiments. Parameters not listed are the generator defaults of
    \cref{sec:exp-workloads}. ``as E1'': \method{full-id}, \method{full-morse},
    \method{gudhi}, \method{drb-id}, \method{drb-id-rb}, \method{mmhm},
    \method{mmhm-norb}; ``DRB methods'': the last four. \method{gudhi} processes
    all updates on instances with at most \NumGudhiFullLimit{} simplices and the
    first $T'={}$\NumPrefix{} updates on larger ones.}
  \label{tab:sweeps}
  \small
  \begin{tabularx}{\linewidth}{@{}l >{\raggedright\arraybackslash}p{3.2cm} >{\raggedright\arraybackslash}X >{\raggedright\arraybackslash}p{4.3cm}@{}}
    \toprule
    & Workload & Sweep & Methods \\
    \midrule
    E0 & samples; all runs of E1--E7; one instance per variant for \method{lcr} & --- & all exact methods; \method{lcr} untimed \\
    E1 & W2, $s=8$ & $N\in\{7,10,15,22,33\}$ & as E1 \\
    E2 & W2, $N=22$ & $s\in\{1,2,4,8,16,32,64\}$ & DRB methods, \method{full-id}, \method{full-morse} \\
    E3 & W3, $k=2$ & $N\in\{14,20,27\}$ & as E1 \\
    E4 & W4 & square, $D=2$: $n\in\{1000,4000,16000\}$; square and cube, $D=3$: $n\in\{1000,4000\}$ & as E1 \\
    E5 & W1 & sphere $L\in\{4,5,6,7\}$; torus $n=m\in\{64,128,256,500\}$ & as E1, and \method{euler} \\
    E6 & W5a, W5b & (a) $L=10^3,3\cdot10^3,10^4,\dots,10^6$; (b) levels $L=3,\dots,8$ & DRB methods, \method{full-id}, \method{gudhi}; \method{euler} on (b) \\
    E7 & W2 $N=22$, W3 $N=20$, W4 square $n=16000$, $D=2$ & $\alpha\in\{0.25,0.5,1,2,4,\infty\}$ \mbox{$\times$} $m\in\{\infty,256,64\}$ & \method{mmhm}, \method{drb-id-rb} \\
    E8 & W2 $N=22$, W3 $N=20$, W4 as in E7, W1 torus $256\times256$ & runs of E1, E3, E4, E5 & DRB methods \\
    \bottomrule
  \end{tabularx}
\end{table}

\subsection{Timing protocol}\label{sec:exp-timing}

The input sequence is parsed completely before timing starts. For the engine
methods, the time of update $t$ is the wall-clock time
(\texttt{std::chrono::steady\_clock}) of processing its block of atomic edits:
the store updates, the work of the method for each edit, the recomputation of
the static baselines, and any rebuild, which is charged to the update that
triggers it. Verification and output are excluded. Building the structure for
$K_0$ is timed separately and reported as build time. For \method{gudhi}, the
time of an update covers the simplex-tree edits, \texttt{compute\_persistence}
and \texttt{betti\_numbers}, again without parsing.

All runs are single-threaded processes. Each configuration ran with
\NumSeeds{} seeds, which change the generated workload; W5a is determined by $L$
alone, so its three runs are repetitions of one instance. Every run processed
$T={}$\NumUpdates{} updates, except \method{gudhi} on instances with more than
\NumGudhiFullLimit{} simplices, which processed the first $T'={}$\NumPrefix{}
(hollow markers in the figures, $'$ in the tables). For each instance and method
we pool the per-update latencies of the three runs and report their median,
interquartile range and 95th and 99th percentiles (p95, p99), and the
\emph{amortised latency}: the total processing time of all updates, rebuilds
included, divided by their number. A \emph{speedup} is a ratio of amortised
latencies; for \method{gudhi} on a prefix, its own prefix is used. Slopes are
least-squares fits of $\log$ latency against $\log$ size. All
\NumConfiguredRuns{} configured runs finished (\NumNotRun{} were not run); they
took \NumComputeHours{} hours of single-threaded processing in total.

\paragraph{Noise.}
Before the sweep, \NumNoiseReps{} identical runs of \method{mmhm} and
\method{full-id} on W2 with $N=22$ gave a coefficient of variation of the median
per-update latency of \NumNoiseCVMmhmMed\,\% and \NumNoiseCVFullIdMed\,\%
(amortised: \NumNoiseCVMmhmAmort\,\% and \NumNoiseCVFullIdAmort\,\%). One
concurrent timed run changed these medians by $\NumNoiseMmhmTwoTimedPct\,\%$ and
$\NumNoiseFullIdTwoTimedPct\,\%$, and three concurrent busy processes by
$\NumNoiseMmhmFourBusyPct\,\%$ and $\NumNoiseFullIdFourBusyPct\,\%$, which is
within the run-to-run spread. At most two timed runs therefore executed at once,
and untimed runs (\method{lcr}, oracle replays) used the remaining two cores.
Differences between methods smaller than the run-to-run spread should not be
read as differences between the methods.

Besides time, every run records per update the number of column additions
(dynamic additions and those of rebuilds separately), the footprints
$|\occ(\sigma)|$ of the deleted simplices, and $\nnz(R)+\nnz(V)$. These are the
quantities of \cref{prop:cost}, and they do not depend on the machine.

\subsection{Correctness protocol}\label{sec:exp-correctness}

Three oracles that share no reduction code check the Betti numbers: (1) the
engine's \method{full-id}, a static recomputation that shares the parser and the
complex store with the dynamic methods but none of their reduction code;
(2) GUDHI, an independent library in another language, which replays the same
sequence; and (3) a pure-Python implementation of the dynamic reduced basis and
of the Morse order, written from the specification and sharing no code with the
engine, together with a dense $\Z$ elimination for small complexes. For W1,
W2, W3 and W5 the closed-form Betti numbers give a fourth check.

Every exact run records the Betti numbers it reports after every update.
\method{full-id} processes all $T$ updates on every instance, and every other
exact run is compared state by state with the \method{full-id} output of the same
instance. The runs are also compared with the closed-form Betti numbers where
these are known (W1, W2, W3, W5), and with the GUDHI oracle on all updates of
instances with at most \NumGudhiFullLimit{} simplices and on the first $T'$
updates of larger ones. In addition, every exact engine run compares itself with
an internal recomputation after $K_0$, after every $50$th update and after the
last one, outside the timed region. On the \NumSamples{} committed samples, all
exact engine configurations (including aggressive rebuild settings), the
reference implementation, GUDHI and the closed-form values are cross-checked on
every state. A mismatch of an exact method stops the whole pipeline.
\method{lcr} runs untimed on one instance of each workload variant and is
compared with \method{full-id} after every update; its \emph{error rate} is the
fraction of updates after which at least one of its Betti numbers is wrong.

Below the level of outputs, the engine's randomised tests check (I1)--(I3) and
compare every exact method with a brute-force rank computation on random edit
sequences, and the reference implementation checks (I1)--(I3) after every
atomic edit, \cref{lem:apparent,cor:critical-work,cor:morse-reduction} on random
complexes and \cref{prop:nonlocal} on the W5 families; it also produced the
counterexamples of \cref{sec:erratum}.


% ---------------------------------------------------------------------------
% Results. One subsection per question of experiments/PLAN.md. Every number is
% a \Num macro (generated/numbers.tex) or sits in a generated table.
% ---------------------------------------------------------------------------

\subsection{E0: Are the exact methods exact, and how wrong is version~1's rule?}\label{sec:e0}

Every exact method reported correct Betti numbers in every state that was
checked. The update rule of version~1 is wrong after almost every update on
W1--W4.

\Cref{tab:correctness} lists the checks. The \NumExactRuns{} runs of exact
methods reported \NumReportedStates{} Betti vectors. \NumVerifiedStates{} of them
were compared with the \method{full-id} output of the same instance,
\NumVerifiedStatesExpected{} with the closed-form values and
\NumVerifiedStatesGudhi{} with GUDHI, and the engine made \NumInternalChecks{}
internal verifications; no check found a mismatch (\NumExactMismatches{}). On
the samples, \NumSampleChecks{} cross-checks between all exact engine
configurations, the reference implementation, GUDHI and the closed-form values
found \NumSampleMismatches{} mismatches.

\method{lcr} reports wrong Betti numbers after \NumLcrErrAll\,\% of all updates
of the E0 instances, and on \NumLcrWorkloadsWrong{} of the \NumLcrWorkloads{}
instances. On W1--W4 it is wrong after
\NumLcrWOneTorusOneHundredTwentyEightErrPct\,\% to
\NumLcrWTwoNFifteenSEightErrPct\,\% of the updates of every instance, for the
first time no later than update \NumLcrWFourPlaneNFourThousandDTwoFirstWrong.
The Betti numbers of W2 are constant, so the errors there come from the algebra
alone. On W5b, \method{lcr} is wrong after \NumLcrWFiveSphereLFiveErrPct\,\% of
the updates: after every deletion of the toggled triangle and after no
reinsertion. The zero column that represents the $2$-cycle of the sphere
absorbed the column of the toggled triangle when $K_0$ was reduced;
\method{lcr} never revisits it, so after each deletion it still reports the
destroyed $2$-cycle, which exists again after each reinsertion. On W5a,
\method{lcr} is never wrong (\NumLcrWFiveCycleLTenThousandErrPct\,\%), and W5a
cannot expose defect~E1. The cycle graph has no free face, so the collapse first
removes the edge with the smallest vertex tuple, $\{0,1\}$, as a critical cell,
which gives it the largest key; W5a toggles exactly this edge, and every
reinsertion gives it the largest key again. No other column ever absorbs the
column with the largest key, so dropping it, as the rule of version~1 does,
leaves a correct reduction. Defect~E1 needs a deleted column that another
reduced column has absorbed, and the updates of W1--W4 delete such columns.

\begin{table}[t]
  \centering
  \caption{Correctness (E0). Upper part: per exact method, the runs, the Betti
    vectors they reported, the states compared with the \method{full-id} output
    of the same instance, with the closed-form values and with GUDHI, and the
    mismatches. Lower part: \method{lcr} (the rule of version~1) on one instance
    of each workload variant, three runs each: updates, updates after which it
    was wrong, and the first wrong update (minimum over the runs).}
  \label{tab:correctness}
  \small
  \generatedtable{tab_correctness}
\end{table}

\subsection{E1: Does the per-update cost depend on the size of the complex at a fixed edit size?}\label{sec:e1}

Only weakly for the DRB methods, and linearly for static recomputation. On W2
with $s=8$, the median per-update latency of the four DRB methods grows with
$|K_0|$ with slope \NumSlopeEOneWTwoSEightMmhmMed{} to
\NumSlopeEOneWTwoSEightDrbIdMed{} over five sizes, against
\NumSlopeEOneWTwoSEightFullIdMed{} for \method{full-id},
\NumSlopeEOneWTwoSEightGudhiMed{} for \method{gudhi} and
\NumSlopeEOneWTwoSEightFullMorseMed{} for \method{full-morse}
(\cref{fig:latency-size}(a)).

At the largest size, $|K_0|={}$\NumWTwoNThirtyThreeSEightSize, \method{mmhm}
takes \NumWTwoNThirtyThreeSEightMmhmAmort\,\textmu s per update amortised,
\NumWTwoNThirtyThreeSEightMmhmSpeedupFullId{} times less than \method{full-id}
(\NumWTwoNThirtyThreeSEightFullIdAmort\,\textmu s) and
\NumWTwoNThirtyThreeSEightMmhmSpeedupGudhi{} times less than \method{gudhi} (on
its prefix); its p99 is \NumWTwoNThirtyThreeSEightMmhmPNinetyNine\,\textmu s,
against \NumWTwoNThirtyThreeSEightFullIdPNinetyNine\,\textmu s for
\method{full-id} (\cref{tab:summary}). The four DRB methods are
indistinguishable on W2: the growth trigger never fired
(\NumWTwoNThirtyThreeSEightMmhmRebuilds{} rebuilds), so \method{mmhm} performs
the same additions as \method{mmhm-norb}, and \method{drb-id-rb} the same as
\method{drb-id}. Building the structure for $K_0$, including the insertion into
the store, took \NumWTwoNThirtyThreeSEightMmhmInitMs\,ms for \method{mmhm} and
\NumWTwoNThirtyThreeSEightFullIdInitMs\,ms for \method{full-id}.

One run is an outlier. At $N=22$, one of the three \method{mmhm-norb} runs was
slower throughout than the other two, although it performed exactly the same
column additions as \method{mmhm} on the same sequence. We attribute this to the
machine, not to the method. It inflates the pooled amortised and p99 latencies
of \method{mmhm-norb} at this size to \NumWTwoNTwentyTwoSEightMmhmNorbAmort{} and
\NumWTwoNTwentyTwoSEightMmhmNorbPNinetyNine\,\textmu s, against
\NumWTwoNTwentyTwoSEightMmhmAmort{} and
\NumWTwoNTwentyTwoSEightMmhmPNinetyNine\,\textmu s for \method{mmhm}.

\paragraph{Static recomputation in the Morse order.}
\method{full-morse} is slower than \method{full-id} on every instance on which
both ran, by factors from \NumWOneSphereLFourFullMorseOverFullId{} (W1 sphere,
$L=4$) to \NumWTwoNThirtyThreeSEightFullMorseOverFullId{} (W2, $N=33$; amortised,
\cref{tab:full-amort}). By \cref{cor:critical-work} its reduction performs
additions only in critical columns, but the collapse after every update costs
more than it saves. For one-shot static recomputation, Morse preprocessing does
not pay off on our workloads; this concerns the ``coreduction-only'' baseline of
version~1 as well.

\begin{figure}[t]
  \centering
  \resultfigure[width=0.92\linewidth]{F2_latency_vs_size}{5.5cm}
  \caption{Median per-update latency against $|K_0|$ at a fixed edit size
    (log--log; band: interquartile range; three runs pooled). (a)~W2 with $s=8$
    (E1); (b)~W3 (E3); (c)~W4 in the square with $D=2$ (E4); (d)~W1 punctured
    sphere (E5). Hollow markers: \method{gudhi} on the prefix $T'$.}
  \label{fig:latency-size}
\end{figure}

\begin{table}[tp]
  \centering
  \caption{Summary at the largest size of each family of E1, E3, E4 and E5:
    $|K_0|$, the number $T$ of timed updates per run, the median, p95 and
    amortised per-update latency in \textmu s (three runs pooled), and the
    speedup over \method{full-id} (ratio of amortised latencies). Every instance
    is listed in \cref{tab:full-amort,tab:full-med,tab:full-p99}.}
  \label{tab:summary}
  \small
  \generatedtable{tab_summary}
\end{table}

\subsection{E2: How does the cost scale with the edit size?}\label{sec:e2}

Linearly for the DRB methods. On W2 with $N=22$ ($|K_0|$ between
\NumWTwoNTwentyTwoSOneSize{} and \NumWTwoNTwentyTwoSSixtyFourSize), the median
per-update latency of the DRB methods grows with the number $s$ of stellar moves
with slope \NumSlopeETwoWTwoNTwentyTwoDrbIdRbMed{} to
\NumSlopeETwoWTwoNTwentyTwoMmhmMed, whereas \method{full-id} and
\method{full-morse} barely depend on $s$ (slopes
\NumSlopeETwoWTwoNTwentyTwoFullIdMed{} and
\NumSlopeETwoWTwoNTwentyTwoFullMorseMed; \cref{fig:latency-s}(a)). The amortised
cost per atomic edit of the DRB methods therefore stays in a narrow band, while
that of recomputation falls roughly as $1/s$ (\cref{fig:latency-s}(b)). There is
no crossover in the measured range: at $s=64$, that is $1024$ atomic edits per
update, the DRB methods are still \NumWTwoNTwentyTwoSSixtyFourMmhmSpeedupFullId{}
to \NumWTwoNTwentyTwoSSixtyFourDrbIdRbSpeedupFullId{} times faster than
\method{full-id}.

\begin{figure}[t]
  \centering
  \resultfigure[scale=1]{F3_latency_vs_edit_size}{4.5cm}
  \caption{Cost against the number $s$ of stellar moves per update (E2; W2,
    $N=22$, $16s$ atomic edits per update; log--log). (a)~Median per-update
    latency with interquartile range; (b)~amortised latency per atomic edit.}
  \label{fig:latency-s}
\end{figure}

\subsection{E3: What does an update cost when the three-dimensional topology changes?}\label{sec:e3}

\method{mmhm} has the lowest amortised latency at every size of W3, but the DRB
is not uniformly better than recomputation: on the smallest instance two DRB
methods are slower than \method{full-id} on average, and every DRB method has a
higher p99.

At $N=14$ ($|K_0|={}$\NumWThreeNFourteenSize), \method{mmhm} takes
\NumWThreeNFourteenMmhmAmort\,\textmu s amortised and \method{full-id}
\NumWThreeNFourteenFullIdAmort\,\textmu s, but \method{drb-id}
(\NumWThreeNFourteenDrbIdAmort\,\textmu s) and \method{mmhm-norb}
(\NumWThreeNFourteenMmhmNorbAmort\,\textmu s) are slower than \method{full-id},
and the p99 of \method{full-id}, \NumWThreeNFourteenFullIdPNinetyNine\,\textmu s,
is below that of every DRB method (from \NumWThreeNFourteenMmhmPNinetyNine\,\textmu s
for \method{mmhm} to \NumWThreeNFourteenDrbIdPNinetyNine\,\textmu s for
\method{drb-id}). At $N=27$ ($|K_0|={}$\NumWThreeNTwentySevenSize) every DRB
method is faster than \method{full-id}, \method{mmhm} by a factor
\NumWThreeNTwentySevenMmhmSpeedupFullId, and every DRB method also has the lower
p99 (\method{mmhm}: \NumWThreeNTwentySevenMmhmPNinetyNine\,\textmu s,
\method{full-id}: \NumWThreeNTwentySevenFullIdPNinetyNine\,\textmu s). At every
size the amortised latencies are ordered \method{mmhm}, \method{drb-id-rb},
\method{mmhm-norb}, \method{drb-id} (\cref{tab:full-amort}): rebuilding matters
more than the order, and the Morse order helps both with and without rebuilds.

W3 is the workload on which the basis grows (\cref{fig:timeseries}(b)). At
$N=20$, without rebuilds, $\nnz(R)+\nnz(V)$ grows over a run from
\NumWThreeNTwentyDrbIdNnzStart{} to \NumWThreeNTwentyDrbIdNnzEnd{}
(\method{drb-id}) and to \NumWThreeNTwentyMmhmNorbNnzEnd{} (\method{mmhm-norb}),
and the latency grows with it. The growth trigger rebuilds \method{mmhm}
\NumWThreeNTwentyMmhmRebuilds{} times and \method{drb-id-rb}
\NumWThreeNTwentyDrbIdRbRebuilds{} times per run (means over the runs), and the
runs end at \NumWThreeNTwentyMmhmNnzEnd{} and \NumWThreeNTwentyDrbIdRbNnzEnd.

\begin{figure}[t]
  \centering
  \resultfigure[width=0.88\linewidth]{F4_time_series}{6cm}
  \caption{Behaviour over time of the four DRB methods (seed~0). Top: per-update
    latency, mean of a sliding window of 50 updates, with the rebuilds of
    \method{mmhm} and \method{drb-id-rb} marked. Bottom: $\nnz(R)+\nnz(V)$ after
    each update. (a)~W2 with $N=22$, $s=8$ (E1); (b)~W3 with $N=20$ (E3);
    (c)~W1 torus $256\times256$ (E5), with the updates after which no hole is
    open marked.}
  \label{fig:timeseries}
\end{figure}

\subsection{E4: How do the methods behave on non-manifold complexes in which every Betti number varies?}\label{sec:e4}

The DRB methods are faster than \method{full-id} on every W4 instance, their
latency hardly depends on $|K_0|$, and the Morse order helps. In the square with
$D=2$ (\cref{fig:latency-size}(c)), their median latency grows with slope
\NumSlopeEFourWFourPlaneDTwoMmhmNorbMed{} to
\NumSlopeEFourWFourPlaneDTwoDrbIdRbMed, against
\NumSlopeEFourWFourPlaneDTwoFullIdMed{} for \method{full-id} (three sizes, so the
slopes are indicative only). Their speedup over \method{full-id} is
\NumWFourSpaceNOneThousandDThreeDrbIdRbSpeedupFullId{} to
\NumWFourPlaneNOneThousandDTwoMmhmNorbSpeedupFullId{} at $n=1000$ and
\NumWFourPlaneNFourThousandDThreeDrbIdSpeedupFullId{} to
\NumWFourPlaneNSixteenThousandDTwoMmhmNorbSpeedupFullId{} at $n\ge4000$. Their
p99 is below that of \method{full-id} on every W4 instance except the two with
$n=1000$ and $D=3$, where it is similar or higher (\cref{tab:full-p99}).

On every W4 instance \method{mmhm-norb} is at least as fast as \method{drb-id}.
At $n=16000$ the amortised latencies are
\NumWFourPlaneNSixteenThousandDTwoMmhmNorbAmort{} and
\NumWFourPlaneNSixteenThousandDTwoDrbIdAmort\,\textmu s, the p99 latencies
\NumWFourPlaneNSixteenThousandDTwoMmhmNorbPNinetyNine{} and
\NumWFourPlaneNSixteenThousandDTwoDrbIdPNinetyNine\,\textmu s, and the p99 of the
column additions per update
\NumEEightWFourPlaneNSixteenThousandDTwoMmhmNorbAddsPNinetyNine{} and
\NumEEightWFourPlaneNSixteenThousandDTwoDrbIdAddsPNinetyNine{} (E8). Under the
default policy, rebuilds fired in only one W4 run of each rebuilding method
(in the cube, $n=1000$, $D=3$). Without rebuilds,
\method{mmhm} does the same additions as \method{mmhm-norb}
(\cref{tab:e8}) but is slower at $n=16000$
(\NumWFourPlaneNSixteenThousandDTwoMmhmAmort{} against
\NumWFourPlaneNSixteenThousandDTwoMmhmNorbAmort\,\textmu s), because it
maintains the index of cells by vertex tuple that a Morse rebuild needs, which
\method{mmhm-norb} switches off.

As anticipated in \cref{sec:exp-workloads}, $\beta_D$ dominates:
at INIT, $\beta_D/n$ is \NumWFourPlaneNFourThousandDThreeBetaDRatio{} to
\NumWFourPlaneNOneThousandDTwoBetaDRatio{} in the square and
\NumWFourSpaceNFourThousandDThreeBetaDRatio{} to
\NumWFourSpaceNOneThousandDThreeBetaDRatio{} in the cube, far above $\beta_0$
and $\beta_1$ (\cref{tab:w4-betti}). These results concern the algebra on
non-manifold complexes, not the topology of point clouds.

\subsection{E5: On surfaces, what does linear algebra cost compared with combinatorial invariants?}\label{sec:e5}

More than recomputation in some cases, and always far more than the
combinatorial counters. On the largest sphere and torus, the insertion-order DRB
methods are slower than \method{full-id} on average; the Morse-ordered ones are
faster, but on no W1 instance by more than a factor
\NumWOneSphereLFiveMmhmSpeedupFullId; every DRB method has a higher p99 than
\method{full-id} on every W1 instance; and \method{euler} takes tens of
microseconds.

On the sphere with $L=7$ ($|K_0|={}$\NumWOneSphereLSevenSize), \method{full-id}
takes \NumWOneSphereLSevenFullIdAmort\,\textmu s amortised, \method{drb-id}
\NumWOneSphereLSevenDrbIdAmort{} and \method{drb-id-rb}
\NumWOneSphereLSevenDrbIdRbAmort\,\textmu s; \method{mmhm-norb} and
\method{mmhm} are \NumWOneSphereLSevenMmhmNorbSpeedupFullId{} and
\NumWOneSphereLSevenMmhmSpeedupFullId{} times faster than \method{full-id}. On
the torus with $n=m=500$ ($|K_0|={}$\NumWOneTorusFiveHundredSize), the speedups
are \NumWOneTorusFiveHundredDrbIdSpeedupFullId{} (\method{drb-id}),
\NumWOneTorusFiveHundredDrbIdRbSpeedupFullId{} (\method{drb-id-rb}),
\NumWOneTorusFiveHundredMmhmSpeedupFullId{} (\method{mmhm}) and
\NumWOneTorusFiveHundredMmhmNorbSpeedupFullId{} (\method{mmhm-norb}). On some
smaller instances even \method{mmhm-norb} is slower than \method{full-id}, for
instance on the sphere with $L=5$ (speedup
\NumWOneSphereLFiveMmhmNorbSpeedupFullId). At $L=7$ the p99 of \method{full-id}
is \NumWOneSphereLSevenFullIdPNinetyNine\,\textmu s and that of \method{mmhm}
\NumWOneSphereLSevenMmhmPNinetyNine\,\textmu s (\cref{tab:full-p99} for all
instances). \method{euler} needs \NumWOneSphereLFourEulerAmort{} to
\NumWOneTorusFiveHundredEulerAmort\,\textmu s per update, including the edits of
the complex store, \NumWOneSphereLSevenEulerSpeedupFullId{} times less than
\method{full-id} on the largest sphere.

The median latency of the DRB methods grows with $|K_0|$ almost like that of
\method{full-id} (slopes \NumSlopeEFiveWOneSphereMmhmNorbMed{} to
\NumSlopeEFiveWOneSphereDrbIdRbMed{} on the sphere, against
\NumSlopeEFiveWOneSphereFullIdMed; \cref{fig:latency-size}(d)): the cost follows
the fundamental $2$-cycle, not the edit. Whenever all holes are closed, which
happens after \NumWOneSphereLSevenClosedStates{} of the \NumUpdates{} updates of
a run on the largest sphere (mean over the runs), the closed surface reappears;
by \cref{prop:nonlocal} the insertion that closes the last hole and the deletion
that opens the next one each cost $\Omega(n_2)$. Between rebuilds the basis also
becomes denser (\cref{fig:timeseries}(c)): on the largest torus,
$\nnz(R)+\nnz(V)$ grows over a run by a factor
\NumWOneTorusFiveHundredDrbIdNnzGrowth{} for \method{drb-id} and
\NumWOneTorusFiveHundredMmhmNorbNnzGrowth{} for \method{mmhm-norb}. This is
\cref{rem:surfaces} in practice: on connected surfaces with known connectivity,
the counters are correct and much cheaper, and the algebra is not needed.

\subsection{E6: Where does locality fail?}\label{sec:e6}

Where \cref{prop:nonlocal} says it does. On both witness families, the cost of
every DRB method per toggle grows linearly with the size of the toggled cycle.
The DRB methods stay faster than \method{full-id} on average, but on the cycle
graph their p99 is higher; only \method{euler} (on the sphere) is nearly
independent of the size.

The amortised latency per toggle grows with slope
\NumSlopeESixWFiveCycleDrbIdAmort{} to \NumSlopeESixWFiveCycleDrbIdRbAmort{}
on W5a ($L$ from $10^3$ to $10^6$) and \NumSlopeESixWFiveSphereDrbIdRbAmort{} to
\NumSlopeESixWFiveSphereMmhmAmort{} on W5b (levels 3 to 8), against
\NumSlopeESixWFiveCycleFullIdAmort{} and \NumSlopeESixWFiveSphereFullIdAmort{}
for \method{full-id} (\cref{fig:nonlocality}). At $L=10^6$, \method{drb-id}
takes \NumWFiveCycleLOneMillionDrbIdAmort\,\textmu s, \method{mmhm}
\NumWFiveCycleLOneMillionMmhmAmort\,\textmu s and \method{full-id}
\NumWFiveCycleLOneMillionFullIdAmort\,\textmu s; on the largest sphere
($|K_0|={}$\NumWFiveSphereLEightSize) the DRB methods take
\NumWFiveSphereLEightDrbIdAmort{} to \NumWFiveSphereLEightMmhmNorbAmort\,\textmu s
and \method{full-id} \NumWFiveSphereLEightFullIdAmort\,\textmu s. On W5a the p99
of every DRB method exceeds that of \method{full-id} at every $L$
(\cref{tab:full-p99}); at $L=10^6$ it is
\NumWFiveCycleLOneMillionMmhmPNinetyNine\,\textmu s for \method{mmhm} against
\NumWFiveCycleLOneMillionFullIdPNinetyNine\,\textmu s. The Morse order does not
help here (at $L=10^6$, \method{drb-id} is faster than \method{mmhm}): the
complex has a single cycle, which every basis must store. \method{euler} takes
\NumWFiveSphereLThreeEulerAmort{} to \NumWFiveSphereLEightEulerAmort\,\textmu s
per toggle on W5b (slope \NumSlopeESixWFiveSphereEulerAmort).

\begin{figure}[t]
  \centering
  \resultfigure[scale=1]{F6_nonlocality}{4.5cm}
  \caption{Amortised latency per toggle on the non-locality witnesses (E6;
    log--log, three runs pooled; the legend gives each least-squares slope).
    (a)~W5a against the cycle length $L$; (b)~W5b against the number $n_2$ of
    triangles. Hollow markers: \method{gudhi} on the prefix $T'$.}
  \label{fig:nonlocality}
\end{figure}

\subsection{E7: How sensitive is \texorpdfstring{\method{mmhm}}{mmhm} to the rebuild policy?}\label{sec:e7}

The answer depends on the workload, and no setting is best everywhere. Where
the basis grows (W3), rebuilding by the growth trigger pays off and the Morse
order at a rebuild helps. Where it does not grow (W2, W4), the trigger hardly
ever fires and periodic rebuilds only add cost. On the large surfaces, rebuilds
make the following updates more expensive (\cref{fig:rebuild},
\cref{tab:e7,tab:rebuild-effect}).

\paragraph{W2 and W4.} With the growth trigger alone, no rebuild fired on W2 at
any $\alpha$, and on W4 one run of three rebuilt once, at $\alpha=0.25$
(\method{mmhm}; \NumESevenWFourPlaneNSixteenThousandDTwoMmhmAQuarterPOffRebuilds{}
rebuilds per run). Periodic rebuilds raised the amortised latency of \method{mmhm} on W2 from
\NumESevenWTwoNTwentyTwoSEightMmhmAOnePOffAmort\,\textmu s ($\alpha=1$, no
period) to \NumESevenWTwoNTwentyTwoSEightMmhmAOnePTwoHundredFiftySixAmort\,\textmu s
($m=256$) and \NumESevenWTwoNTwentyTwoSEightMmhmAOnePSixtyFourAmort\,\textmu s
($m=64$). A Morse rebuild, which includes a collapse, costs more than an
insertion-order one: \method{drb-id-rb} took
\NumESevenWTwoNTwentyTwoSEightDrbIdRbAOnePSixtyFourAmort\,\textmu s at $m=64$.

\paragraph{W3.} At $\alpha=1$ without a period, \method{mmhm} took
\NumESevenWThreeNTwentyMmhmAOnePOffAmort\,\textmu s amortised against
\NumESevenWThreeNTwentyMmhmAOffPOffAmort\,\textmu s without rebuilds, and
\method{drb-id-rb} \NumESevenWThreeNTwentyDrbIdRbAOnePOffAmort{} against
\NumESevenWThreeNTwentyDrbIdRbAOffPOffAmort\,\textmu s; the p99 of
\method{mmhm} fell from \NumESevenWThreeNTwentyMmhmAOffPOffPNinetyNine{} to
\NumESevenWThreeNTwentyMmhmAOnePOffPNinetyNine\,\textmu s. Too small a threshold
hurts again: at $\alpha=0.25$ \method{mmhm} rebuilt
\NumESevenWThreeNTwentyMmhmAQuarterPOffRebuilds{} times per run and its p99 was
\NumESevenWThreeNTwentyMmhmAQuarterPOffPNinetyNine\,\textmu s. \method{mmhm} has
a lower amortised latency than \method{drb-id-rb} at every $(\alpha,m)$ of the
grid. In the 100 updates after a rebuild, the rebuilding method performs fewer
additions than its never-rebuilding twin on the same updates (ratios
\NumRebuildEffectWThreeNTwentyMmhmWindow{} for \method{mmhm} and
\NumRebuildEffectWThreeNTwentyDrbIdRbWindow{} for \method{drb-id-rb} at $N=20$).

\paragraph{W1.} Under the default policy rebuilds fire on every W1 instance
(\cref{tab:rebuild-effect}). They lower the amortised latency on the smaller
instances and raise it on the largest: \method{mmhm} is faster than
\method{mmhm-norb} on the sphere with $L\le5$ and the tori with $n\le256$ (for
instance \NumWOneSphereLFiveMmhmAmort{} against
\NumWOneSphereLFiveMmhmNorbAmort\,\textmu s at $L=5$) but slower on the sphere
with $L\ge6$, and about equal on the torus with $n=500$; \method{drb-id-rb} is faster than \method{drb-id} on the sphere
with $L\le6$ and the tori with $n\le128$, and slower on the others. On the sphere
with $L\ge6$ and the tori with $n\ge256$, the 100 updates after an \method{mmhm}
rebuild perform more additions than the same updates of \method{mmhm-norb}, by
factors from \NumRebuildEffectWOneSphereLSixMmhmWindow{} to
\NumRebuildEffectWOneTorusFiveHundredMmhmWindow; for \method{drb-id-rb} this
happens on the large spheres (\NumRebuildEffectWOneSphereLSixDrbIdRbWindow{} and
\NumRebuildEffectWOneSphereLSevenDrbIdRbWindow) but not on the tori. We have not
identified the mechanism.

\begin{figure}[t]
  \centering
  \resultfigure[width=0.88\linewidth]{F7_rebuild_sensitivity}{6cm}
  \caption{Rebuild policy (E7): amortised (top) and p99 (bottom) per-update
    latency of \method{mmhm} and \method{drb-id-rb} against the growth threshold
    $\alpha$, one line per period $m$; three runs pooled. (a)~W2 with $N=22$,
    $s=8$; (b)~W3 with $N=20$; (c)~W4 in the square with $n=16000$, $D=2$. With
    $\alpha=m=\infty$ (``off'') the methods never rebuild. The grid, with the
    rebuild counts, is in \cref{tab:e7}.}
  \label{fig:rebuild}
\end{figure}

\subsection{E8: What does the Morse order buy?}\label{sec:e8}

Lighter tails of the column additions on every workload measured, and smaller
median footprints on W4 and W2. Not every quantile improves, and with rebuilds
the insertion-order rebuild leaves the lighter footprint tails. We compare the two methods that never rebuild, \method{mmhm-norb} (Morse
order) and \method{drb-id} (insertion order), on the non-rebuild updates of the
E1, E3, E4 and E5 runs (\cref{fig:footprint}, \cref{tab:e8}); the footprint of an
update is the sum of $|\occ(\sigma)|$ over its deleted simplices.

On W4 ($n=16000$, $D=2$), the Morse order lowers the p99 of the additions per
update from \NumEEightWFourPlaneNSixteenThousandDTwoDrbIdAddsPNinetyNine{} to
\NumEEightWFourPlaneNSixteenThousandDTwoMmhmNorbAddsPNinetyNine{} and the median
footprint from \NumEEightWFourPlaneNSixteenThousandDTwoDrbIdFootprintMed{} to
\NumEEightWFourPlaneNSixteenThousandDTwoMmhmNorbFootprintMed. On W2 ($N=22$) it
lowers the p99 of the additions from
\NumEEightWTwoNTwentyTwoSEightDrbIdAddsPNinetyNine{} to
\NumEEightWTwoNTwentyTwoSEightMmhmNorbAddsPNinetyNine{} and the median footprint
from \NumEEightWTwoNTwentyTwoSEightDrbIdFootprintMed{} to
\NumEEightWTwoNTwentyTwoSEightMmhmNorbFootprintMed, but raises the median
additions from \NumEEightWTwoNTwentyTwoSEightDrbIdAddsMed{} to
\NumEEightWTwoNTwentyTwoSEightMmhmNorbAddsMed{} and the p99 footprint from
\NumEEightWTwoNTwentyTwoSEightDrbIdFootprintPNinetyNine{} to
\NumEEightWTwoNTwentyTwoSEightMmhmNorbFootprintPNinetyNine. On W3 ($N=20$) the
median additions fall from \NumEEightWThreeNTwentyDrbIdAddsMed{} to
\NumEEightWThreeNTwentyMmhmNorbAddsMed, and on the W1 torus ($256\times256$)
the p99 of the additions from
\NumEEightWOneTorusTwoHundredFiftySixDrbIdAddsPNinetyNine{} to
\NumEEightWOneTorusTwoHundredFiftySixMmhmNorbAddsPNinetyNine; on both the median
footprints are almost equal. With rebuilds the picture changes: on W3 and on the
W1 torus, the insertion-order rebuild of \method{drb-id-rb} leaves lighter
footprint tails than the Morse rebuild of \method{mmhm} (p99
\NumEEightWThreeNTwentyDrbIdRbFootprintPNinetyNine{} against
\NumEEightWThreeNTwentyMmhmFootprintPNinetyNine{} on W3), although \method{mmhm}
had the lower amortised latency on W3 (E7). \Cref{lem:apparent} and its
corollaries describe the state right after a rebuild; between rebuilds the Morse order is a
heuristic, and these measurements are its whole justification.

\begin{figure}[t]
  \centering
  \resultfigure[width=0.88\linewidth]{F5_footprint_ccdf}{6cm}
  \caption{What the Morse order buys (E8): complementary cumulative
    distributions of the footprint per update (top; the sum of
    $|\occ(\sigma)|$ over its deletions) and of the column additions per update
    (bottom), over the updates without a rebuild (three runs pooled; log--log).
    (a)~W2 with $N=22$, $s=8$; (b)~W3 with $N=20$; (c)~W4 in the square with
    $n=16000$, $D=2$. On W2 and W4 no rebuild fired, so \method{mmhm} coincides
    with \method{mmhm-norb} and \method{drb-id-rb} with \method{drb-id}.}
  \label{fig:footprint}
\end{figure}
